\section{Introduction}
\label{sec:introduction}

\subsection{The Mixing Time Problem}

Markov chain Monte Carlo (MCMC) redistricting has become a standard tool
for detecting partisan gerrymandering.
The inference framework is simple: draw a large ensemble of redistricting
plans from the uniform distribution over valid plans, compute a partisan
metric for each plan, and report where the enacted map falls in the
distribution.
A map in the 95th percentile for Republican advantage is likely gerrymandered.

The statistical validity of this inference depends entirely on whether the
chain has \emph{mixed}---whether the empirical distribution from the chain
approximates the true uniform distribution over valid plans.
If the chain has not mixed, the ensemble does not represent the uniform
distribution; it represents the neighborhood of the starting plan.

This paper asks: how long must a redistricting chain run to mix?
We give two complementary answers: a theoretical bound from Markov chain
spectral theory, and empirical measurements from \recom{} chains on
actual census-tract graphs.

\subsection{Main Results}

\paragraph{Theoretical bound.}
Census-tract adjacency graphs are approximately planar (they are derived
from geographic tract boundaries).
For random walks on planar graphs, the spectral gap satisfies
$\sgap = \Omega(1/n^2)$~\citep{lipton1979,levin2009}, where $n$ is the
number of vertices (tracts).
This implies mixing time $\tmix = O(n^2 \log n)$.
For a state with $n = 2{,}000$ tracts, the theoretical bound is
$\tmix = O(4 \times 10^6 \log 2{,}000) \approx O(3 \times 10^7)$ steps.

\paragraph{Empirical result.}
In practice, \recom{} chains achieve $\rhat < 1.1$ (the G.4 convergence
criterion) in 2{,}000--20{,}000 steps.
The gap between theory ($10^7$) and practice ($10^4$) is explained by the
concentration of the stationary distribution near high-compactness plans:
the chain rarely needs to visit the vast low-compactness regions of state
space, so practical mixing is much faster than the worst-case bound.

\paragraph{ConvergenceSweep alternative.}
The \cs{} algorithm~\citep{b16paper} provides a deterministic $O(T \cdot n)$
stopping criterion for plan \emph{generation}.
For plan \emph{evaluation} (auditing enacted maps), ensemble methods remain
irreplaceable.
We identify the boundary between these two regimes.

\subsection{Relation to Prior Work}

\citet{deford2021} introduced \recom{} and studied its stationary
distribution theoretically, but did not provide explicit mixing time bounds.
\citet{autry2021} studied the forest-based recombination variant and proved
polynomial mixing for certain graph families, but the bounds are not tight
for planar graphs.
\citet{chikina2017} showed that some tests are valid without full mixing,
but their conditions are hard to verify in practice.
This paper provides the first explicit spectral-gap-based mixing time
lower bounds for the redistricting MCMC context.

\subsection{Paper Structure}

Section~\ref{sec:mcmc} formally defines the \recom{} Markov chain and
its state space.
Section~\ref{sec:bounds} derives the $\Omega(1/n^2)$ spectral gap bound
and the resulting mixing time.
Section~\ref{sec:empirical} presents empirical mixing measurements for
six states.
Section~\ref{sec:sweep} contrasts MCMC mixing with \cs{}.
Section~\ref{sec:when-ensemble} identifies the irreplaceable ensemble use cases.
Section~\ref{sec:conclusion} concludes.
